% FILE: 01_research_question.tex
% Project: sources-wasm4pm-compat
% Role: type-foundry-and-witness-lattice-authority

\section{Research Question}
\label{sec:sources-wasm4pm-compat-question}

\subsection{Problem Statement}
\label{sec:sources-wasm4pm-compat-question-problem}

Process execution engines face a structural problem at their admission layer: they cannot,
by construction, prevent the intake of structurally invalid evidence without either
(a)~encoding admissibility rules as runtime checks that can be bypassed or
(b)~rejecting all evidence wholesale. The conventional approach---runtime validation with
conditional error returns---leaves the critical invariant that ``evidence is admissible''
as a dynamic property, not a type-level guarantee. This means that a caller can construct
an \texttt{Evidence} block, skip the validation step, and pass it directly to a downstream
function without a compile error.

The problem is compounded at the boundary between a compatibility layer and an execution
engine. When two components are separated by a WebAssembly component boundary, the only
guarantees that survive the crossing are those encoded in the WIT interface type system.
If the Rust type system has enforced an admissibility invariant, but the WIT boundary
definition does not reflect that invariant, the guarantee evaporates at the crossing.

The research question is therefore: \emph{Can the admissibility boundaries of process
evidence be encoded as compile-time type laws that are non-forgeable, receipted, and
preserved across WebAssembly component boundaries?}

\subsection{Old-Regime Assumption Challenged}
\label{sec:sources-wasm4pm-compat-question-oldregime}

The assumption being challenged is that admissibility validation belongs in the execution
engine's runtime path. Under the old regime, an execution engine receives a raw event log,
runs a suite of validation functions, and either proceeds or returns an error. The validation
is decoupled from the type of the data: the same \texttt{EventLog} type is passed to the
engine regardless of whether it has been validated. There is no structural distinction
between a validated log and an unvalidated one.

This assumption is operationally expedient but epistemologically weak. Under the Van der
Aalst constitution---if the code says it worked but the event log cannot prove a lawful
process happened, then it did not work---a runtime validation path that is not itself
receipted is not evidence. It is an unwitnessed claim.

\texttt{wasm4pm-compat} challenges this assumption by asserting that the type system is
the validation layer, that compile-fail fixtures are the receipts, and that the WIT
boundary definition is the contract that makes the guarantee portable across component
boundaries.

\subsection{Hypothesis}
\label{sec:sources-wasm4pm-compat-question-hypothesis}

The hypothesis is threefold. First, a triadic cryptographic container
\texttt{Evidence<T, State, Witness>} binding a payload, a process-model state, and a
lattice-monotonic witness via BLAKE3 hash and Ed25519 signature can serve as the
sole admissibility unit for process evidence, making non-admissible evidence
inexpressible at the type level. Second, a bounded join-semilattice \texttt{(W, \(\sqsubseteq\),
\(\sqcup\), \(\bot\), \(\top\))} over the eight \texttt{WitnessMarker} algebraic proof
structures can enforce replay soundness (Axiom~2) such that a non-monotonic witness
transition is rejected at the type boundary, not at runtime. Third, these guarantees
can be materialized as WebAssembly Interface Type boundary definitions that survive
component model crossings intact.

The hypothesis is AUTHOR THESIS: it is grounded in the architecture described in
\texttt{type-law-atlas.md} (v30.1.1) and \texttt{research-verdict.md} (v30.1.2), but
full integration with the \texttt{wasm4pm} execution engine remains a downstream obligation
as of 2026-06-01.

\subsection{Success Criteria}
\label{sec:sources-wasm4pm-compat-question-criteria}

The success criteria as defined in \texttt{research-verdict.md} (v30.1.2, 2026-05-31)
across five audit categories are:

\begin{enumerate}
  \item \textbf{Witness Lattice Completeness}: The algebraic structure is sound across
    all five modeled process family domains (Petri Nets, BPMN~2.0, POWL~2.0, Process
    Trees, Declare constraints). Status: \textbf{PASS}.
  \item \textbf{Evidence Triadic Container Admissibility}: The container is cryptographically
    non-forgeable and enforces valid state transitions. Status: \textbf{PASS}.
  \item \textbf{Admission/Refusal Law Enforcement}: Boundary control operates under
    default-deny; all refusal pathways are documented and logged. Status: \textbf{PASS}.
  \item \textbf{Loss Policy Thermodynamics}: Permissible and absolute loss boundaries
    are mathematically defined with self-halt semantics on terminal loss. Status: \textbf{PASS}.
  \item \textbf{Non-Forgeability Guarantees}: Cryptographic binding (Axiom~1), lattice
    monotonicity (Axiom~2), and signature admissibility (Axiom~3) are fully specified.
    Status: \textbf{COMPLETE}.
\end{enumerate}

The ALIVE gate is defined by \texttt{config.toml} as 406 compile-pass fixtures and
398 compile-fail fixtures (804 total). Two structural defects remain open---the WfNet
split-brain and zero multi-step pipeline witness-consistency fixtures---placing the
overall project status at \textbf{PARTIAL} rather than ALIVE.
